\subsection{P010 --- BIM-to-BRICK: Using graph modeling for IoT/BMS and spatial semantic data interoperability within digital data models of buildings}
\label{paper:P010}
\textbf{Authors:} Filippo Vittori, Chuan Fu Tan, Anna Laura Pisello, Adrian Chong, Cristina Piselli, Clayton Miller\par
\textbf{Major Category:} BIM and Semantic Information\par
\textbf{Secondary Category:} BMS, BAS, BACS and Building Automation, Sensors, IoT and Data Integration, Digital Twin\par
\textbf{Keywords:} BIM, Brick Schema, BMS, IoT, Knowledge graph, Data interoperability, Digital Twin\par
\paragraph{主な主張}
BIM-to-BRICKはBIM情報をknowledge graphへ自動変換し, 事例では932 BIM instancesと17 subjectsの実験dataを458 BRICK objects, 1219 relationshipsへ17 sで統合し, BIMとBRICK間のbidirectional linkを構築した. \cite{Vittori:2025BIMtoBRICK}
\paragraph{新規性}
BIMとBMS/IoTのsemantic interoperabilityに必要なmappingをgraph modelingにより自動化し, manual interventionを削減するBIM-to-BRICK workflowを提示した点に特徴がある.
\paragraph{位置付け}
BIMからBrick knowledge graphへのsemantic変換とinteroperabilityが中核であるためBIM and Semantic Informationを主分類とする.
